\begin{textsample}{POS Dim 4 – 2019\_c5 – Score 1.00 – 2019\_c5/enade\_2019\_farm\_a84c9952ee.txt}  \label{ex:f4_pos_193}
CAPÍTULO 5 DISTRIBUIÇÃO DOS CONCEITOS 5.1 PANORAMA NACIONAL DA DISTRIBUIÇÃO DOS CONCEITOS O Conceito Enade é um indicador de qualidade que avalia os cursos por intermédio dos desempenhos dos estudantes a partir dos resultados do Enade.
Os valores possíveis vão de 1 ( pior situação ) a 5 ( melhor situação ).
Os cursos com apenas um ou sem nenhum concluinte participante não obtêm o Conceito Enade, ficando Sem Conceito ( SC ).
Na Tabela 5.1, são apresentadas a quantidade e a distribuição de cursos de Farmácia participantes do Enade/2019, por faixa de conceito e Grande Região.
Enfatiza -se, mais uma vez, que a diferença entre os cursos tabulados neste capítulo e no capítulo 2 corresponde aos cursos que não foram avaliados, ou seja, àqueles, em princípio, sem alunos concluintes que tivessem sido inscritos para a prova.
Observando -se os dados da Tabela 5.1, nota -se que, dos 431 cursos participantes, 178 ( 41,3 \% ) classificaram -se com conceito 3, o valor modal.
Esse foi também o conceito modal nas regiões Nordeste ( 39,8 \% ), Sudeste ( 46,3 \% ) e Centro-Oeste ( 44,4 \% ).
Nas demais regiões, a classe modal foi a seguinte ; Norte conceito 2 ( 52,9 \% ) e Sul conceito 4 ( 40,3 \% ).
O conceito 4 foi o segundo mais frequente em nível nacional ( 27,6 \%, correspondendo a 119 cursos ), e o conceito 2, o terceiro ( 21,6 \%, correspondendo a 93 cursos ).
Houve, ainda, 28 cursos ( 6,5 \% ) que receberam conceito 5, e 11 cursos ( 2,6 \% ) que receberam conceito 1.
Dos 431 cursos de Farmácia, dois ( 0,5 \% ) ficaram Sem Conceito ( SC ).
A região Norte participou com 34 cursos ou 7,9 \% do total nacional.
Desses, 18 cursos receberam o conceito 2, o que equivale a 52,9 \% do total regional, além de ser o conceito modal na região, como já comentado.
A dois cursos ( 5,9 \% ) atribuiu -se conceito 1, e os conceitos 3 e 4 foram atribuídos a 7 cursos cada ( 20,6 \% para cada ).
Nenhum curso ficou Sem Conceito ( SC ) e nenhum curso recebeu o conceito 5 nessa região.
A região Nordeste participou com 83 cursos ou 19,3 \% do total nacional.
Nessa região, 33 cursos ( 39,8 \% em termos regionais ) obtiveram conceito 3, o conceito modal para a região, como já comentado.
Os conceitos 2 e 4 foram atribuídos, respectivamente, a 25 e a 16 cursos ( 30,1 \% e 19,3 \%, respectivamente ).
Ainda nessa região, oito cursos ( 9,6 \% ) foram avaliados com o conceito 5, e um recebeu conceito 1 ( 1,2 \% ).
Nenhum curso ficou Sem Conceito ( SC ).
Dos 188 cursos participantes da região Sudeste ( 43,6 \% do total nacional ), o conceito 3 foi o conceito modal, como já anteriormente assinalado, tendo sido obtido por 87 cursos ( 46,3 \% ).
O conceito 1 foi atribuído a quatro cursos ( 2,1 \% ), e o conceito 2, a 25 cursos ( 13,3 \% ).
Receberam os conceitos 4 e 5, respectivamente, 58 cursos ( 30,9 \% ) e 12 cursos ( 6,4 \% ).
Nessa região, dois cursos ficaram Sem Conceito ( SC ).
Os 72 cursos da região Sul corresponderam a 16,7 \% do total nacional.
Como já indicado, houve predominância do conceito 4, atribuído a 40,3 \% dos cursos da região ( conceito modal ), ou seja, atribuído a 29 dos 72 cursos participantes na região Sul.
O conceito 1 foi atribuído a um curso ( 1,3 \% ), o conceito 2 a nove cursos ( 12,5 \% ), o conceito 3 a 27 cursos ( 37,5 \% ) e o conceito 5, a seis cursos ( 8,3 \% cada ).
Nenhum curso ficou Sem Conceito ( SC ) na região Sul.
Já dos 54 cursos participantes na região Centro-Oeste ( 12,5 \% do total nacional ), 24 ( 44,4 \% em termos regionais ) receberam conceito 3, o conceito modal, como já fora destacado.
Três cursos ( 5,6 \% ) obtiveram conceito 1, 16 cursos ( 29,6 \% ) obtiveram conceito 2, nove ( 16,7 \% ), conceito 4 e outros dois ( 3,7 \% ) receberam conceito 5.
Nenhum curso ficou Sem Conceito ( SC ).
No Gráfico 5.1, é apresentada a distribuição cumulativa do Conceito Enade segundo Grande Região.
Em um gráfico de distribuição cumulativa, em cada valor de abscissa ( eixo x ), o valor apresentado no eixo das ordenadas ( eixo y ) é o valor acumulado da variável ( conceito Enade ).
Observando a linha vertical pontilhada, que ajuda a visualizar o percentual de conceitos acumulados até o valor 3, ou seja, conceitos 1, 2 e 3, pode -se notar que, por exemplo, as regiões Norte ( linha azul ) e Centro-Oeste ( linha laranja ) apresentam aproximadamente 80 \% dos cursos com conceito 3 ou abaixo deste valor.
A região Sul ( linha rosa ), por outro lado, apresenta 51 \% dos cursos na mesma situação : com conceito 3 ou abaixo.
Destaca -se que o conceito 5 não foi alcançado por nenhum curso na região Norte ( linha azul ), na qual todos os cursos têm conceito menor ou igual a 4 e, por isso, a poligonal cumulativa alcança 100 \% no conceito 4.
Em linhas gerais, em um gráfico de distribuição cumulativa, poligonais mais à esquerda representam Grandes Regiões com pior distribuição, e poligonais mais à direita, Grandes Regiões com melhores conceitos.
A região Norte ( linha azul ) apresenta os cursos com conceitos mais baixos, ao passo que as regiões Sudeste ( linha vermelha ) e Sul ( linha rosa ) apresentam as distribuições com valores maiores.
Cursos em IES situadas nas regiões Nordeste ( linha verde ) e Centro-Oeste ( linha laranja ) apresentam uma situação intermediária. 5.2 CONCEITOS POR CATEGORIA ADMINISTRATIVA, POR MODALIDADE DE ENSINO E GRANDE REGIÃO Na Tabela 5.2, é apresentada a distribuição dos cursos de Farmácia participantes do Enade/2019, por Categoria Administrativa e Modalidade de ensino, de acordo com os conceitos por eles alcançados, segundo as Grandes Regiões brasileiras.
Dos 431 cursos participantes, 70 ( 16,2 \% ) eram ministrados em Instituições Públicas, e 361 ( 83,8 \% ), em Instituições Privadas.
Quanto à Modalidade de ensino, todos os cursos são presenciais.
De acordo com as informações da Tabela 5.2, em termos nacionais, entre Instituições Públicas e Privadas, observa -se que, dos 28 cursos avaliados com conceito 5, 22 eram oferecidos em IES Públicas, e seis, em IES Privadas.
Dos 70 cursos participantes de IES Públicas, o conceito 4 foi o valor modal, atribuído a 39 cursos ( 55,7 \% ).
Dos demais cursos de Instituições Públicas participantes, dois cursos foram avaliados com conceito 2, sete cursos ( 10,0 \% ) foram avaliados com conceito 3 e, como já comentado, 22 cursos ( 31,4 \% ), com conceito 5.
Nenhum curso ficou Sem Conceito ( SC ).
Nesta categoria, nenhum curso recebeu conceito 1.
Na rede privada, o conceito modal foi 3, atribuído a 171 cursos, o correspondente a 47,4 \% dos 361 cursos da categoria.
Dos demais cursos participantes, 11 ( 3,0 \% ) receberam conceito 1, e 91 ( 25,2 \% ), conceito 2.
O conceito 4 foi atribuído a 80 cursos ( 22,2 \% ), e o conceito 5, a seis cursos ( 1,7 \% ).
Nesta Categoria Administrativa, dois cursos ( 0,6 \% ) ficaram Sem Conceito ( SC ).
Dos 431 cursos participantes da modalidade presencial, o conceito 3 foi o valor modal, atribuído a 178 cursos ( 41,3 \% ).
Dos demais cursos nesta modalidade de ensino, 11 cursos ( 2,6 \% ) receberam conceito 1, 93 cursos ( 21,6 \% ) receberam conceito 2, 119 cursos ( 27,6 \% ), conceito 4, e 28 cursos ( 6,5 \% ), conceito 5.
Nesta modalidade de ensino, dois cursos ( 0,5 \% ) ficaram Sem Conceito ( SC ).
Na análise por região, observa -se que, na região Norte, as Instituições Privadas participaram com 27 cursos ( 79,4 \% do total regional ).
O conceito modal para as Instituições Privadas na região foi 2, atribuído a 17 cursos, correspondendo a 63,0 \% dessa categoria na região.
Dois cursos foram avaliados com conceito 1, sete cursos com conceito 3, e um curso, com conceito 4.
Nenhum curso ficou Sem Conceito ( SC ) ou recebeu conceito 5.
As Instituições Públicas participaram com sete cursos na região Norte ( 20,6 \% do total regional ), aos quais foram atribuídos os conceitos 2 e 4.
Nenhum curso ficou Sem Conceito ( SC ) ou com conceito 1, 3 ou 5 nessa combinação de categoria e região.
Todos os cursos oferecidos na região Norte eram presenciais.
Na região Nordeste, a rede privada concentrou 65 dos 83 cursos participantes, o equivalente a 78,3 \% do total da região.
Dentre os cursos oferecidos por IES Privadas no Nordeste, 33 cursos foram avaliados com conceito 3, o conceito modal.
Um curso obteve conceito 1, 25 cursos obtiveram conceito 2, cinco cursos receberam conceito 4 e um curso ficou com conceito 5.
Nenhum curso ficou Sem Conceito ( SC ).
As Instituições Públicas da região Nordeste participaram com 18 cursos ( 21,7 \% do total da região ).
Desses, 11 obtiveram conceito 4, o valor modal, seguidos por outros sete cursos, que receberam conceito 5.
Nenhum curso ficou Sem Conceito ( SC ) ou com os demais conceitos nessa região.
Todos os cursos oferecidos na região Nordeste eram presenciais.
Na região Sudeste, a proporção de cursos da rede privada, 88,8 \%, correspondendo a 167 dos 188 cursos participantes.
Nessa categoria e região, o conceito modal foi 3, atribuído a 84 cursos, e dois cursos ficaram Sem Conceito ( SC ).
Os demais foram avaliados com conceito 1 ( quatro cursos ), conceito 2 ( 25 cursos ), conceito 4 ( 48 cursos ) e conceito 5 ( quatro cursos ).
Dos 21 cursos oferecidos por Instituições Públicas na região Sudeste, o conceito modal foi 4, atribuído a dez cursos.
Nessa região, nenhum curso de IES Públicas ficou Sem Conceito ( SC ), três cursos receberam conceito 4, e oito, conceito 5.
Todos os cursos oferecidos na região Sudeste eram presenciais.
As Instituições Privadas concentraram 58 dos 72 cursos participantes da região Sul, 80,6 \% do total regional.
Desses, 25 cursos obtiveram conceito 3, o conceito modal.
Nesta combinação de Categoria Administrativa e Grande Região, um curso recebeu conceito 1, nove cursos receberam conceito 2, outros 22, conceito 4, e um, conceito 5.
Nenhum curso ficou Sem Conceito ( SC ).
As Instituições Públicas na região Sul participaram com 14 cursos ( 19,4 \% dos cursos da região ), aos quais foram atribuídos os conceitos 3, 4 ou 5.
Nenhum curso da região Sul oferecido por IES Pública ficou Sem Conceito ( SC ) ou recebeu os conceitos 1 e 2.
Todos os cursos oferecidos na região Sul eram presenciais.
Na região Centro-Oeste, 44 dos 54 cursos participantes eram de Instituições Privadas ( 81,5 \% em termos regionais ).
Desses, 22 receberam conceito 3, o conceito modal.
Dos demais cursos, três receberam conceito 1, 15 receberam conceito 2, e outros quatro, conceito 4.
Nenhum curso ficou Sem Conceito ( SC ), e nenhum curso recebeu conceito 5.
Quanto aos dez cursos oferecidos por Instituições Públicas na região Centro-Oeste ( 18,5 \% do total regional ), cinco cursos apresentaram conceito 4, conceito modal, um recebeu conceito 2, dois receberam conceito 3 e dois receberam conceito 5.
Nenhum curso ficou Sem Conceito ( SC ) ou recebeu os demais conceitos nessa combinação de categoria e região.
No Gráfico 5.2, apresenta -se a distribuição do Conceito Enade, segundo a Categoria Administrativa da IES.
Os cursos em IES Públicas ( linha azul ) apresentam uma poligonal mais à direita do que os em IES Privadas ( linha verde ) e, consequentemente, uma distribuição de Conceitos Enade com valores maiores. 5.3 CONCEITOS POR ORGANIZAÇÃO ACADÊMICA E GRANDE REGIÃO Na Tabela 5.3, encontra -se a distribuição dos conceitos atribuídos aos cursos participantes do Enade/2019, na Área de Farmácia, por Organização Acadêmica, segundo as Grandes Regiões brasileiras.
Dos 431 cursos de Farmácia participantes, 186 eram oferecidos em Universidades, 103 em Centros Universitários, 140 em Faculdades e 2 em CEFET/IF.
Esta distribuição corresponde a, respectivamente, 43,2 \%, 23,9 \%, 32,5 \% e 0,5 \% dos cursos.
De acordo com os dados apresentados, de todos os 28 cursos avaliados com conceito 5, 25 eram vinculados a Universidades.
Os cursos vinculados a esse tipo de Organização Acadêmica tiveram o conceito 4 como conceito modal, atribuído a 72 cursos ( 38,7 \% ).
Os demais cursos vinculados a Universidades receberam os conceitos 1 ( um curso ), 2 ( 17 cursos ), 3 ( 70 cursos ) e conceito 5 ( 25 cursos, como já mencionado ).
Ficou Sem Conceito ( SC ), um curso.
Dos cursos em Centros Universitários, o conceito modal foi 3, atribuído a 45 cursos ( 43,7 \% ).
Nenhum curso vinculado a esse tipo de Organização Acadêmica ficou Sem Conceito ( SC ), e os outros cursos receberam os conceitos 1 ( um curso ), 2 ( 27 cursos ), 4 ( 29 cursos ) e 5 ( um curso ).
Dos 140 cursos mantidos por Faculdades, 62 ( 44,3 \% ) receberam conceito 3, o conceito modal.
O segundo conceito mais frequentemente atribuído aos cursos oferecidos por essa categoria de IES foi o conceito 2, recebido por 49 cursos.
Na sequência, 18 cursos obtiveram conceito 4, nove cursos, conceito 1, e um curso, conceito 5.
Dos cursos oferecidos em Faculdades, um ficou Sem Conceito ( SC ).
Dos 2 cursos mantidos por CEFET/IF, um recebeu conceito 3 e o outro recebeu conceito 5.
Considerando -se, separadamente, as regiões brasileiras, verifica -se que, na região Norte, as Universidades participaram com dez dos 34 cursos avaliados, aos quais foi atribuído conceito 2 para quatro cursos e conceito 4 para seis cursos, que é a classe modal.
Nenhum curso ficou Sem Conceito ( SC ) ou recebeu os demais conceitos nesta combinação de categoria e região.
Os Centros Universitários da região Norte foram representados por sete cursos, os quais obtiveram conceitos 2, 3 e 4.
Nenhum curso ficou Sem Conceito ( SC ) ou recebeu os demais conceitos nesta combinação de categoria e região.
As Faculdades participaram com 17 cursos na região Norte, os quais receberam os conceitos 1 ( dois cursos ), 2 ( dez cursos, conceito modal ) e 3 ( cinco cursos ).
Nenhum curso ficou Sem Conceito ( SC ), e nenhum curso oferecido em Faculdades da região obteve conceito 4 e 5.
Na região Nordeste, as Universidades participaram com 23 dos 83 cursos da Área de Farmácia oferecidos.
O conceito modal 4 foi atribuído a 11 cursos.
Os cursos restantes obtiveram conceitos 3 ( quatro cursos ) e 5 ( oito cursos ).
Nenhum curso oferecido por Universidades no Nordeste ficou Sem Conceito ( SC ) ou recebeu conceito 1 ou 2.
Os Centros Universitários contaram com 22 cursos participantes na região Nordeste, 11 dos quais receberam o conceito modal 3.
Para os demais cursos, foram atribuídos os conceitos 2 ( oito cursos ) e 4 ( três cursos ).
Nenhum curso oferecido por Centros Universitários no Nordeste ficou Sem Conceito ( SC ) ou recebeu os conceitos 1 ou 5.
As Faculdades foram representadas por 38 cursos na região Nordeste, 18 dos quais receberam conceito 3, o conceito modal.
Para os demais cursos, foram atribuídos os conceitos 1 ( um curso ), 2 ( 17 cursos ) e 4 ( dois cursos ).
Na região Nordeste, nenhum curso mantido por Faculdades ficou Sem Conceito ( SC ) ou recebeu conceito 5.
Na região Sudeste, as Universidades concentraram 89 dos 188 cursos de Farmácia da região.
Dos cursos oferecidos em Universidades, nessa região, o conceito modal foi 3, atribuído a 42 cursos.
Os demais cursos receberam os conceitos 2 ( oito cursos ), 4 ( 29 cursos ) e 5 ( nove cursos ).
Um curso ficou Sem Conceito ( SC ) e nenhum recebeu conceito 1.
Os Centros Universitários participaram com 46 cursos na região Sudeste, dos quais 21 obtiveram o conceito modal 3.
O conceito 2 foi atribuído a nove cursos, o conceito 4, a 15 cursos, e o conceito 5, a um curso.
Nenhum curso ficou Sem Conceito ( SC ) ou recebeu conceito 1 nessa combinação de organização acadêmica e região.
Já as Faculdades foram representadas por 52 cursos na região Sudeste, dos quais um ficou Sem Conceito ( SC ).
O conceito modal 3 foi atribuído a 24 cursos.
Aos demais cursos foram atribuídos os conceitos 1 ( quatro cursos ), 2 ( oito cursos ), 4 ( 14 cursos ) e 5 ( um curso ).
Os CEFET/IF participaram com um curso, recebendo o conceito 5.
Dos 72 cursos da região Sul, 46 eram mantidos por Universidades, dos quais o conceito modal foi 4, atribuído a 20 cursos.
Os demais cursos receberam os conceitos 1 ( um curso ), 2 ( quatro cursos ), 4 ( 15 cursos ) e 5 ( seis cursos ).
Nesse tipo de organização, nenhum curso da região Sul ficou Sem Conceito ( SC ).
Dos 16 cursos participantes de Centros Universitários da região Sul, nenhum ficou Sem Conceito ( SC ), e nenhum recebeu os conceitos 1 ou 5.
Os conceitos restantes, 3 e 4, foram atribuídos a sete cursos cada, e o conceito 2 foi a atribuído a dois cursos.
Quanto aos 9 cursos vinculados a Faculdades na região Sul, quatro receberam o conceito modal 3.
O conceito 2 foi atribuído a três cursos, e o conceito 4, a dois cursos.
Nenhum curso ficou Sem Conceito ( SC ) ou recebeu conceito 1 ou 5.
Já os CEFET/IF participaram com um curso na região Sul, ao qual foi atribuído o conceito 3.
Na região Centro-Oeste, 18 dos 54 cursos oferecidos eram mantidos por Universidades.
Para este tipo de Organização Acadêmica, nenhum curso ficou Sem Conceito ( SC ) ou recebeu conceito 1.
Os conceitos 2 e 4 foram atribuídos a, respectivamente, um e seis cursos, e o conceito 3, conceito modal, foi atribuído a nove cursos.
O conceito 5 foi atribuído a dois cursos.
Os Centros Universitários da região Centro-Oeste contaram com 12 cursos : os conceitos 2 e 3 foram atribuídos a quatro cursos cada.
Para os demais cursos, foram atribuídos os conceitos 1 ( um curso ) e 4 ( três cursos ).
Nesta região, nenhum curso oferecido em Faculdades ficou Sem Conceito ( SC ), e nenhum obteve conceito 5.
Quanto aos 24 cursos vinculados a Faculdades na região Centro-Oeste, 11 receberam igualmente os conceitos 2 e 3.
O conceito 1 foi atribuído a um curso.
Nenhum curso ficou Sem Conceito ( SC ) ou recebeu conceito 4 ou 5.
Na região Centro-Oeste, nenhum curso foi vinculado a CEFET/IF.
No Gráfico 5.3, apresenta -se a distribuição do Conceito Enade dos cursos de Farmácia, segundo a Organização Acadêmica da IES.
Os cursos em CEFET/IF ( linha laranja ) apresentam uma poligonal mais à direita do que as demais, denotando uma melhor distribuição de Conceitos Enade.
No outro extremo, com distribuição muito semelhante entre si, encontram -se os cursos mantidos por Universidades ( linha azul ), Centros Universitários ( linha verde ) e Faculdades ( linha vermelha ), apresentando as respectivas poligonais mais à esquerda e, consequentemente, distribuição de Conceitos Enade com menores valores.

% matched lemmas: 
\end{textsample}
